\section{Discussion and Limitations}

\label{sec:discussion}

The controlled results establish two things clearly. First, the deployed hybrid selector is consistently and substantially outperformed by BM25 at every token budget, with differences of -36.1111 percentage points, -29.3651 percentage points, and -33.7302 percentage points, all surviving Holm correction. Second, removing only the temporal term recovers most of this deficit, localizing the cause. These two facts together motivate CTLS: the temporal term is the primary source of the ranking error, and a fix that targets it precisely rather than discarding it globally is the right architectural response.

The mechanism analysis turns the empirical observation into a formal condition. A fixed query-independent temporal weight can outweigh a positive Jaccard advantage whenever the freshness gap exceeds a threshold determined by the weight ratio and the half-life. This condition is satisfied routinely in single-session factual extraction because target facts are old and distractors are recent. The evidence recall deficit mirrors the accuracy deficit exactly: the displaced items are the relevant older ones, and once they leave the prompt the answering model cannot produce the correct answer. CTLS addresses this by partitioning candidates before temporal scoring, guaranteeing that every positive-overlap item is considered for inclusion before any zero-overlap item.

The relationship between CTLS and hybrid\_no\_time deserves careful statement. hybrid\_no\_time solves temporal dominance by eliminating the temporal term entirely, which is effective but discards useful information. CTLS solves the same problem more precisely by restricting temporal scoring to within-tier comparisons: a fresher item can still outrank a less-fresh item when both have positive overlap, but it cannot outrank an item from a different tier with lower overlap. For queries where all candidates have positive overlap, CTLS and hybrid\_no\_time are equivalent and both outperform the deployed hybrid. For mixed-tier queries -- the more common case -- CTLS is strictly stronger. The hybrid\_no\_time result is therefore a conservative lower bound on what CTLS will achieve when directly evaluated.

The reuse experiment shows that construction cost amortizes across queries. The total token ratio of the model-acknowledgement path to the raw path is 1.21 at the endpoint after 160 distinct queries. CTLS does not change this curve: it is a selector-only substitution with zero construction overhead. An application that serves many queries against a single memory instance will pay the write cost once and amortize it; CTLS provides better retrieval quality at every query at no additional cost.

Several restrictions apply. The controlled evaluation uses a frozen single-layer bank; the production rail has different capacity constraints. The judge substitution and stratified subset are disclosed deviations from the official evaluation protocol. The blinded human-adjudication task remains incomplete and results use the primary predeclared judge. CTLS is established theoretically with empirical support from the hybrid\_no\_time ablation; a direct CTLS implementation and controlled evaluation remains the primary recommended follow-up. The present artifact records candidate scores and skipped-item behavior and can directly support that experiment.
